\section{Skills: módulos de capacidad reutilizables}

\subsection{¿Qué es un Skill?}

Un Skill es, en esencia, un \textbf{patrón reutilizable de ejecución de tareas}. Boris explica que, en la demostración anterior en la que Co-Work generó una hoja de cálculo de Excel, en realidad cargó un skill de Excel preinstalado; esa es la razón por la que el modelo sabía cómo operar Excel de manera correcta.

\begin{importantbox}{La propuesta de valor de un Skill}
Si usas una herramienta o un formato de archivo poco común (como AutoCAD o Salesforce), puedes escribir un Skill personalizado para enseñarle a Claude cómo operarlo. Una instancia de Co-Work con más Skills puede manejar una gama más amplia de tareas y produce resultados de mayor calidad.
\end{importantbox}

\subsection{Recomendaciones de uso}

La recomendación de Boris es \textbf{empezar simple antes de ir a lo complejo}:

\begin{enumerate}
    \item Basta con instalar Co-Work y la extensión de Chrome para empezar a usarlo
    \item No dediques tiempo a personalizar un Skill desde el principio
    \item Considera escribir el Skill correspondiente solo cuando notes que Co-Work no funciona bien en un caso concreto
    \item Los Skills son más adecuados para usuarios avanzados como una manera de optimizar el rendimiento
\end{enumerate}

\subsection{Resumen de la sección}

Un Skill es la capa de capacidad extensible de Co-Work y de Claude Code. Permite que el usuario codifique conocimiento de un dominio en módulos reutilizables, lo que mejora el desempeño del agente en tareas específicas.

